\section{Calibrating the operating-cost ratio \texorpdfstring{$\varrho$}{rho}}
\label{app:rho}

The single price parameter of the adoption block, $\varrho=\bar P^E_G/\bar P^E_B$,
is a delivered operating cost: the cost of one unit of energy service (a kilometre)
from the green durable relative to the brown one. It is not a price per unit of
energy (per kilowatt-hour, electricity is dearer than gasoline) but a cost per unit
of service, net of each technology's conversion efficiency. The upfront cost of
switching is a separate object, $\psi_g$.

We calibrate $\varrho$ on the electric-vehicle margin, the durable for which both
the operating-cost gap and the adoption elasticity (Section~\ref{sec:taste_id}) are
measured. With $P^G,P^B$ the retail electricity and gasoline prices and $e_G,e_B$
real-world energy use,
\[
  \varrho=\frac{P^G\,e_G}{P^B\,e_B}
        =\frac{0.290\times0.21}{1.62\times0.070}
        =\frac{0.0609}{0.1134}=0.54 .
\]

\begin{table}[H]\centering\small
\begin{tabular}{llll}
\toprule
 & & Value & Source \\
\midrule
Electricity price, households & $P^G$ & 0.290 EUR/kWh & Eurostat \texttt{nrg\_pc\_204}, 2025S2, incl.\ taxes \\
Gasoline price, Euro-95       & $P^B$ & 1.62 EUR/L    & EC Weekly Oil Bulletin, 2025S2, incl.\ taxes \\
Electric-car energy use       & $e_G$ & 0.21 kWh/km   & IEA \textit{Global EV Outlook 2026} \\
Gasoline-car fuel use         & $e_B$ & 7.0 L/100km   & National on-road data (e.g.\ FR SDES) \\
\bottomrule
\end{tabular}
\caption{Inputs to $\varrho$. European, tax-inclusive, same semester.}
\label{tab:rho_inputs}
\end{table}

Three choices are worth stating. \emph{(i) Tax-inclusive prices.} The baseline has
no carbon tax, so the household compares retail prices; the gap between the two
ratios ($0.54$ inclusive against $0.88$ excluding taxes) is the asymmetric
taxation of the two fuels, which the ETS acts on rather than assumes away.
\emph{(ii) Real-world energy use.} Both figures are on-road or metered rather than
type-approval, so charging losses are already in the electricity figure and the two
sides rest on the same basis. \emph{(iii) Fleet averages.} A matched compact pair
(electric $17$~kWh/100km, gasoline $6.2$~L/100km) gives $0.49$, so the ratio does
not turn on vehicle-size composition. As an internal check, $0.21$~kWh/km is
$2.3$~litres of gasoline equivalent per 100~km, $68\%$ below the gasoline car's
$7.0$, matching the roughly $70\%$ energy saving the same source reports for a
comparable electric car. The heating margin (an air-source heat pump against a
condensing gas boiler at Eurostat household electricity and gas prices) gives a
nearby $\varrho\approx0.62$, so reading the durable as a heat-pump/EV composite
would move the ratio only inside the range below.

\paragraph{Sensitivity.} Holding the steady-state green share at $\bar D^G=5\%$,
$\psi_g$ re-solves as $\varrho$ varies. Across the whole plausible range the policy
CEV ranking is unchanged (flat transfer $\succ$ Slutsky $\succ$ ex-post cap
$\succ$ ex-ante ETS $\succ$ laissez-faire $\succ$ green subsidy); only the crisis
adoption magnitude responds, by about one percentage point of peak green share.

\begin{table}[H]\centering\small
\begin{tabular}{lccccc}
\toprule
$\varrho$                       & 0.48 & 0.54 & 0.62 & 0.70 & 0.80 \\
\midrule
$\psi_g$ (solved)               & 0.827 & 0.770 & 0.696 & 0.625 & 0.538 \\
peak $\Delta D^G$ (pp)          & 7.44  & 7.72  & 8.10  & 8.46  & 8.87 \\
peak output loss $-y$ (\%)      & 2.85  & 2.86  & 2.87  & 2.88  & 2.90 \\
\bottomrule
\end{tabular}
\caption{Sensitivity of the baseline to $\varrho$. Price shock, import booking.
$\bar D^G=5\%$ imposed throughout; the policy ordering (not shown) is invariant.}
\label{tab:rho_sensitivity}
\end{table}
